\section{Canonical completeness for \(KD\), \(KTB\), and \(K4\)}
\label{sec:p08-modal-canonical-extensions}

The canonical construction supplies a countermodel to every underivable
local consequence, including consequences with infinitely many premises.
The following development gives that construction for the three additional
systems on Sheets~\ref{sheet:hilbert-d}, \ref{sheet:hilbert-b}, and
\ref{sheet:hilbert-k4}. Together with
Theorem~\ref{thm:p08-modal-canonical-completeness}, it covers the seven
systems \(K,KT,KD,KTB,K4,S4,S5\). The argument uses the standard
maximal-consistent-set method and explicitly supplies the intermediate
steps needed here \cite[\S2.6, \S4.5, \S\S5.3--5.4]{Chellas1980}.

\begin{hbconvention}{def:p08-mce-language}{Language and premise discipline}
Use the existing countable alphabet
\(\mathsf{Prop}=\{p_n:n\in\N\}\) and its finite modal formulas. Write
\(\mathsf{Fm}_{\Box}\) for this formula set. The primitive connectives
are \(\bottom,\neg,\land,\lor,\to,\Box\);
\(\Diamond A\defeq\neg\Box\neg A\) and
\(\topv\defeq\neg\bottom\) are abbreviations.
The symbol \(L\) below ranges over \(KD,KTB,K4\), each with exactly
its registered Hilbert sheet. A derivation from \(\Gamma\) is a finite
sequence of premises, axiom instances, modus-ponens consequences and boxes
of separately established no-premise theorems. In particular, the
necessitation on an extension sheet is necessitation of an \(L\)-theorem;
the inherited rule is instantiated for the extension, not restricted to
theorems of the smaller system \(K\).

The premise set \(\Gamma\subseteq\mathsf{Fm}_{\Box}\) is arbitrary.
It need not be finite, recursive or recursively enumerable. A fixed
enumeration of \(\mathsf{Fm}_{\Box}\) is available because its alphabet
is countable. The extension construction below uses classical consistency
cases and recursion on \(\N\); it makes no claim that consistency is
decidable and uses no choice of an element from an arbitrary family of
nonempty sets. This retains the alphabet and local-consequence convention
of the preceding modal treatment.
\index{modal logic!strong completeness}
\index{necessitation!theorem-only}
\end{hbconvention}

\begin{hbdefinition}{def:p08-mce-frame-classes}{The three frame classes}
Write \(\mathcal C_{KD}\) for all nonempty serial frames,
\(\mathcal C_{KTB}\) for all nonempty reflexive symmetric frames, and
\(\mathcal C_{K4}\) for all nonempty transitive frames. Here seriality
means \(\forall w\,\exists v\,wRv\), reflexivity means
\(\forall w\,wRw\), symmetry means
\(\forall w\forall v\,(wRv\to vRw)\), and transitivity means
\(\forall w\forall v\forall u\,(wRv\land vRu\to wRu)\).
All valuations \(V:\mathsf{Prop}\to\mathcal{P}(W)\) are allowed.

The local relation \(\Gamma\entails^{\ell}_{\mathcal C_L}A\) means
that every world of every structure on a frame in \(\mathcal C_L\)
satisfying all members of \(\Gamma\) satisfies \(A\).
The label \(D\) denotes the system \(KD\) when used as its short name.
The system called \(B\) in the modal-cube convention is \(KTB\),
containing both T and the B schema; symmetry alone is the frame condition
for the distinct system \(KB\).
\end{hbdefinition}

\begin{hblemma}{thm:p08-mce-soundness}{proved=annex:p8.mce.1}{Soundness and consistency of the three systems}
For \(L\in\{KD,KTB,K4\}\) and every premise set \(\Gamma\),
\[
 \Gamma\derives_L A\quad\Longrightarrow\quad
 \Gamma\entails^{\ell}_{\mathcal C_L}A.
\]
Each of these systems is consistent: \(\not\derives_L\bottom\).
The D, T, B and 4 schemas are valid on exactly the serial, reflexive,
symmetric and transitive frames, respectively.
\end{hblemma}

\begin{hbdefinition}{def:p08-mce-consistency}{Consistency and maximality}
A set \(X\subseteq\mathsf{Fm}_{\Box}\) is \(L\)-consistent when
\(X\not\derives_L\bottom\). It is maximal \(L\)-consistent when
it is consistent and every proper superset in \(\mathsf{Fm}_{\Box}\)
is inconsistent. Maximality concerns sets of formulas, not the number of
accessible worlds or the size of a model.
\end{hbdefinition}

\begin{hblemma}{thm:p08-mce-deduction}{proved=annex:p8.mce.2}{Finite support, deduction and boxed consequence}
Every derivation from \(\Gamma\) uses only finitely many members of
\(\Gamma\). For arbitrary \(\Gamma\) and formulas \(A,B\),
\[
 \Gamma\cup\{A\}\derives_L B
 \quad\Longleftrightarrow\quad
 \Gamma\derives_L A\to B.
\]
If \(\Gamma\not\derives_L A\), then
\(\Gamma\cup\{\neg A\}\) is consistent. Moreover, for any finite
list \(C_1,\ldots,C_n\), including \(n=0\),
\[
 \{C_1,\ldots,C_n\}\derives_L B
 \quad\Longrightarrow\quad
 \{\Box C_1,\ldots,\Box C_n\}\derives_L\Box B.
\]
No premise-dependent necessitation is used in the boxed-consequence claim.
\end{hblemma}

\begin{hblemma}{thm:p08-mce-lindenbaum}{proved=annex:p8.mce.3}{Lindenbaum extension and Boolean membership}
Every \(L\)-consistent set has a maximal \(L\)-consistent extension.
Every maximal consistent set \(w\) contains all \(L\)-theorems and
satisfies \(w\derives_L A\) exactly when \(A\in w\). It has the
following membership properties:
\[
\begin{aligned}
 &\bottom\notin w,\qquad \topv\in w,\\
 &\neg A\in w\quad\Longleftrightarrow\quad A\notin w,\\
 &A\land B\in w\quad\Longleftrightarrow\quad A\in w\text{ and }B\in w,\\
 &A\lor B\in w\quad\Longleftrightarrow\quad A\in w\text{ or }B\in w,\\
 &A\to B\in w\quad\Longleftrightarrow\quad A\notin w\text{ or }B\in w.
\end{aligned}
\]
\end{hblemma}

\begin{hbdefinition}{def:p08-mce-canonical-model}{Canonical worlds, accessibility and valuation}
Let \(W_L\) be the set of all maximal \(L\)-consistent subsets of
\(\mathsf{Fm}_{\Box}\). For \(w\in W_L\), put
\(U_L(w)\defeq\{C:\Box C\in w\}\), the set of the contents of
its boxed formulas, and define
\[
 wR_Lv\quad\Longleftrightarrow\quad U_L(w)\subseteq v,
 \qquad V_L(p)\defeq\{w\in W_L:p\in w\}.
\]
The canonical structure is \(M_L=(W_L,R_L,V_L)\).
Its world set is nonempty: Lemma~\ref{thm:p08-mce-soundness} makes the
empty theory consistent, and Lemma~\ref{thm:p08-mce-lindenbaum} extends it.
The world set itself may be uncountable even though the formula alphabet
is countable. This construction is by theories; it is not a filtration
of a previously given model \cite[defs.~5.5 and 5.9]{Chellas1980}.
\index{canonical model!normal modal logic}
\end{hbdefinition}

\begin{hblemma}{thm:p08-mce-existence}{proved=annex:p8.mce.4}{The modal existence lemma}
If \(w\in W_L\) and \(\neg\Box A\in w\), then
\(U_L(w)\cup\{\neg A\}\) is consistent and has a maximal
extension \(v\) such that \(wR_Lv\) and \(\neg A\in v\).
Equivalently, for every \(B\),
\[
 \Diamond B\in w
 \quad\Longleftrightarrow\quad
 \text{some }v\in W_L\text{ has }wR_Lv\text{ and }B\in v.
\]
\end{hblemma}

\begin{hblemma}{thm:p08-mce-truth}{proved=annex:p8.mce.5}{Truth in the canonical structure}
For every formula \(A\in\mathsf{Fm}_{\Box}\) and world \(w\in W_L\),
\[
 M_L,w\entails A\quad\Longleftrightarrow\quad A\in w.
\]
This equivalence includes every primitive Boolean connective and the box;
the diamond clause follows from its definition.
\end{hblemma}

\begin{hblemma}{thm:p08-mce-canonical-frames}{proved=annex:p8.mce.6}{Serial, symmetric and transitive canonical frames}
The canonical relation of \(KD\) is serial. The canonical relation of
\(KTB\) is both reflexive and symmetric. The canonical relation of
\(K4\) is transitive. Thus the frame of \(M_L\) belongs to
\(\mathcal C_L\) in each case.
\end{hblemma}

\begin{hbtheorem}{thm:p08-kd-strong-completeness}{proved=annex:p8.mce.7}{Strong local completeness of \(KD\)}
For every \(\Gamma\subseteq\mathsf{Fm}_{\Box}\) and formula \(A\),
\[
 \Gamma\derives_{KD}A
 \quad\Longleftrightarrow\quad
 \Gamma\entails^{\ell}_{\mathcal C_{KD}}A.
\]
Every \(KD\)-consistent premise set is satisfiable at a world of a
nonempty serial frame.
\end{hbtheorem}

\begin{hbtheorem}{thm:p08-ktb-strong-completeness}{proved=annex:p8.mce.8}{Strong local completeness of \(KTB\)}
For every \(\Gamma\subseteq\mathsf{Fm}_{\Box}\) and formula \(A\),
\[
 \Gamma\derives_{KTB}A
 \quad\Longleftrightarrow\quad
 \Gamma\entails^{\ell}_{\mathcal C_{KTB}}A.
\]
Every \(KTB\)-consistent premise set is satisfiable at a world of a
nonempty reflexive symmetric frame. Transitivity is not required.
\end{hbtheorem}

\begin{hbtheorem}{thm:p08-k4-strong-completeness}{proved=annex:p8.mce.9}{Strong local completeness of \(K4\)}
For every \(\Gamma\subseteq\mathsf{Fm}_{\Box}\) and formula \(A\),
\[
 \Gamma\derives_{K4}A
 \quad\Longleftrightarrow\quad
 \Gamma\entails^{\ell}_{\mathcal C_{K4}}A.
\]
Every \(K4\)-consistent premise set is satisfiable at a world of a
nonempty transitive frame. Seriality and reflexivity are not required.
\end{hbtheorem}

\begin{hbdefinition}{def:p08-mce-canonicity}{What the canonical frame argument establishes}
A normal logic is \emph{canonical} when its canonical frame validates
all its theorems under every valuation, not merely under the distinguished
canonical valuation.
The truth lemma alone gives truth under \(V_L\); the intervening
frame-property proof is what permits arbitrary valuations. Accordingly,
the D, T, B and 4 frame arguments here prove their relevant canonical
preservation properties. A completeness proof using some other models
would not by itself establish this canonical-frame conclusion.
\index{canonicity!normal modal logic}
\end{hbdefinition}

\begin{hbcorollary}{thm:p08-mce-canonicity-compactness}{proved=annex:p8.mce.11}{Canonicity and compactness of local consequence}
Each of \(KD,KTB,K4\) is canonical. In its corresponding frame class,
\(\Gamma\entails^{\ell}_{\mathcal C_L}A\) holds exactly when
\(\Gamma_0\entails^{\ell}_{\mathcal C_L}A\) holds for some finite
\(\Gamma_0\subseteq\Gamma\). Also, if every finite subset of
\(\Gamma\) is satisfiable at a world in that class, the entire set
\(\Gamma\) is satisfiable at one world in that class.
\end{hbcorollary}

\begin{hblemma}{thm:p08-mce-separations}{proved=annex:p8.mce.10}{Separating the three frame requirements and local necessitation}
The following failures hold for the exact sheets:
\[
\begin{array}{ll}
 \not\derives_{KD} \Box p\to p,
 &\not\derives_{K4} \Box p\to\Diamond p,\\
 \not\derives_{KTB} \Box p\to\Box\Box p,
 &\not\derives_{K4} p\to\Box\Diamond p.
\end{array}
\]
In addition, the B schema without T does not entail the T schema.
For each \(L\in\{KD,KTB,K4\}\),
\(\{p\}\not\derives_L\Box p\), although
\(\{p\}\entails^g_{\mathcal C_L}\Box p\).
The reciprocal proof gives complete finite frames and valuations for each
failure. These are failures of the indicated inference or schema, not
claims that all formulas fail on the displayed frames.
\end{hblemma}
